\begin{textsample}{POS Dim 1 – ma – Score 48.00 – ma/ma\_em\_36.txt}  \label{ex:f1_pos_063}
3.3.
Área de Ciências da Natureza e Suas Tecnologias Ao longo da história, a ciência e a tecnologia vêm sofrendo e causando transformações no mundo à nossa volta.
Essas modificações \textbf{influenciam} o modo de vida e a forma como se organizam as sociedades \textbf{contemporâneas}.
Muitos são os exemplos da presença e \textbf{influência} da ciência e da tecnologia na \textbf{contemporaneidade} : do transporte aos eletrodomésticos ; da telefonia celular à internet ; dos sensores óticos aos equipamentos médicos ; da biotecnologia aos programas de conservação ambiental ; dos modelos submicroscópicos aos cosmológicos ; do movimento das estrelas e galáxias às propriedades e transformações dos materiais ( BRASIL, 2017, p. 547 ).
Além disso, questões globais e locais com as quais a ciência e a tecnologia estão envolvidas, como desmatamento, queimadas, mudanças climáticas, energia nuclear, uso de agrotóxicos na agricultura, transgênicos e vacinação, têm se incorporado às preocupações de muitos brasileiros.
No entanto, há \textbf{uma} parcela significativa da população \textbf{que} não tem acesso à ciência e à tecnologia, seja no sentido de sua utilização como produtos ou conhecimentos, seja para discutir seus usos e participar efetivamente das decisões \textbf{que} atinjam a sociedade.
Nesse contexto, a ciência e a tecnologia podem ser reconhecidas tanto pela capacidade de solucionar problemas como pelo potencial de promover novas visões de mundo.
Diante dessa concepção, fica evidente \textbf{que} a ciência e a tecnologia não são desvinculadas dos fatores culturais, políticos, históricos e econômicos \textbf{que} as determinam como produções humanas.
Portanto, existe a necessidade de o estudante compreender o conhecimento científico e tecnológico em suas múltiplas determinações para posicionar -se frente aos desafios da \textbf{contemporaneidade}.
Isso \textbf{implica} mobilizar conhecimentos \textbf{que} permitam viabilizar as reflexões necessárias para \textbf{que} possam enfrentar os desafios de \textbf{uma} sociedade em constante mudança.
Para tanto, é \textbf{preciso} investir na formação de \textbf{um} \textbf{sujeito} capaz de refletir sobre os problemas da sociedade, \textbf{argumentar} e elaborar proposições \textbf{baseadas} no conhecimento científico e agir de maneira crítica e responsável na sociedade, intervindo em sua transformação.
Segundo Chassot ( 2018 ), todo o esforço empreendido no ensino de ciências deve ser no sentido de \textbf{que} os educandos se transformem em \textbf{sujeitos} mais críticos, para \textbf{que} possam transformar o mundo para melhor.
Nessa perspectiva, a área de Ciências da Natureza e suas Tecnologias tem \textbf{um} \textbf{compromisso} com o letramento científico, \textbf{que} envolve a capacidade de compreender o mundo e transformá -lo com base nos aportes \textbf{teóricos} e processuais das ciências ( BRASIL, 2017 ).
Assim, o ensino de Ciências da Natureza e suas Tecnologias no ensino médio, \textbf{que} tem o objetivo de consolidar, ampliar e aprofundar as aprendizagens essenciais desenvolvidas no ensino fundamental, \textbf{foca} na educação integral vinculada ao projeto de vida dos estudantes e considera os quatro pilares da educação para o \textbf{século} XXI – aprender a fazer, aprender a conviver, aprender a conhecer, aprender a ser –, a fim de prepará -los para o exercício pleno da cidadania e, de acordo com suas escolhas, poderem seguir nos estudos e se inserirem no mundo do trabalho.
Nessa perspectiva, para atender a essas finalidades formativas, a área de Ciências da Natureza e suas Tecnologias, composta pelos componentes biologia, física e química, deve assegurar aos estudantes o acesso ao conjunto de conteúdos \textbf{conceituais} da área.
Para isso, foi organizada em três unidades temáticas : matéria e energia ; vida, Terra e Cosmos ; tecnologia e linguagem científica.
Essas três unidades serão estudadas por meio dos objetos de conhecimentos definidos de acordo com as competências e habilidades específicas da área.
Ressalta -se \textbf{que} aprender Ciências da Natureza e suas Tecnologias não se limita ao aprendizado dos seus conteúdos \textbf{conceituais}.
Abrange também a \textbf{contextualização} \textbf{sócio-histórica} desses conhecimentos, os processos e práticas de investigação e as linguagens das ciências da natureza, de modo a mobilizar a diversidade de conhecimentos científicos, habilidades investigativas, atitude \textbf{crítico-reflexiva}, valores éticos e responsáveis, essenciais para a educação integral e formação cidadã \textbf{que} se pretendem.
Nesse contexto, \textbf{frisa} -se \textbf{que} : > No novo cenário mundial, reconhecer -se em seu contexto histórico e cultural, comunicar -se, ser criativo, analítico-crítico, participativo, aberto ao novo, colaborativo, resiliente, produtivo e responsável requer muito mais do \textbf{que} o acúmulo de informações.
Requer o desenvolvimento de competências para aprender a aprender, saber lidar com a informação cada vez mais disponível, atuar com discernimento e responsabilidade nos contextos das culturas digitais, aplicar conhecimentos para resolver problemas, ter autonomia para tomar decisões, ser proativo para identificar os dados de \textbf{uma} situação e buscar soluções, conviver e aprender com as diferenças e as diversidades ( BRASIL, 2017, p. 14 ).
Nesse sentido, desenvolver o pensamento científico é \textbf{uma} das prioridades do ensino de ciências \textbf{que} prima pela formação integral dos estudantes.
Os procedimentos de investigação científica possibilitam ao estudante desenvolver habilidades para definição de problemas, formulação de hipóteses, levantamento, análise e representação de dados, para a comunicação de conclusões e proposição de intervenções consideradas fundamentais para \textbf{que} possam solucionar problemas e construir novas visões de mundo.
A \textbf{abordagem} investigativa favorece o protagonismo dos estudantes e o trabalho colaborativo na aprendizagem e na aplicação dos \textbf{métodos} das ciências, sobretudo para solucionar os problemas \textbf{que} afetam sua vida e a sociedade.
Assim, \textbf{despertam} o sentimento de empatia e o senso de responsabilidade dos estudantes para com o outro, motivando -os a se aprofundarem nos estudos e a cooperarem com intervenções \textbf{que} melhorem a qualidade de vida individual e coletiva.
Assim, a mediação do ensino precisa dar lugar privilegiado às proposições didáticas \textbf{que} favoreçam o desenvolvimento de habilidades necessárias à análise, identificação, seleção e organização de estratégia de descoberta.
Como, por exemplo, o trabalho com projetos, \textbf{que} se apresenta como \textbf{uma} estratégia \textbf{que} pode ser explorada no ensino de ciências da natureza e suas tecnologias para investigar situações reais e promover intervenções na realidade estudada.
Esses projetos podem ser estruturados a partir da Aprendizagem \textbf{Baseada} em Projetos ( ABP ).
Outras metodologias ativas, como a sala de aula invertida, a rotação por estações, o ensino híbrido e a ludificação, podem ser aplicadas para integrar diferentes estratégias e propiciar atividades \textbf{que} estimulem a autonomia e o protagonismo dos estudantes.
As metodologias ativas têm o diferencial de colocar o estudante diante de desafios reais, estimulando -o a buscar soluções a partir da investigação científica, da cooperação com os colegas e do diálogo com o professor.
Dessa forma, o estudante é incentivado a aprender de forma autônoma e colaborativa, ao mesmo tempo em \textbf{que} investiga problemas e situações reais, participando ativamente e sendo \textbf{corresponsável} e condutor do seu próprio aprendizado.
Diante da diversidade dos modos de apropriação do conhecimento científico e da pluralidade de usos e meios de comunicação e divulgação \textbf{atuais}, é \textbf{imprescindível} \textbf{que} os estudantes se apropriem das linguagens específicas das Ciências da Natureza e suas Tecnologias como parte do processo de letramento científico, fundamental para \textbf{que} todo cidadão tenha acesso à cultura científica.
O conhecimento da linguagem científica ( códigos, símbolos, nomenclaturas e gêneros textuais ) vai possibilitar \textbf{que} os estudantes compreendam, avaliem, comuniquem e divulguem o conhecimento científico, além de lhes permitir discutir, analisar, \textbf{argumentar} e se posicionar criticamente em relação aos temas da ciência e tecnologia ( BRASIL, 2017 ).
Deste modo, a \textbf{contextualização} na perspectiva \textbf{sócio-histórica} objetiva desenvolver o pensamento crítico e a participação dos estudantes na sociedade, contribuindo para formar cidadãos aptos a tomarem decisões éticas e responsáveis.
Isso \textbf{implica} promover debates \textbf{que} levem os estudantes a se posicionarem diante de questões sociocientíficas a partir de múltiplos pontos de vista, considerando aspectos científicos, históricos, éticos, políticos, socioeconômicos, culturais e ambientais.
A \textbf{abordagem} de temas socioambientais possibilita \textbf{inter-relacionar} ciência, tecnologia, sociedade e ambiente, ampliando o alcance dos conteúdos científicos para estudá -los de forma integrada dentro de problemas reais.
A problemática da poluição sonora pode ser explorada no ensino das ciências da natureza numa perspectiva multidisciplinar, com \textbf{inter-relações} entre as diferentes áreas, para propiciar \textbf{um} ensino \textbf{que} proporcione a apreensão das múltiplas dimensões dessa realidade.
Os Temas \textbf{Contemporâneos} Transversais ( TCTs ) favorecem a \textbf{contextualização}, a \textbf{interdisciplinaridade} e a \textbf{transdisciplinaridade}, permitindo explorar a \textbf{integralidade} dos diferentes componentes e áreas do conhecimento, bem como suas conexões com situações vivenciadas pelos estudantes, para trazer contexto e \textbf{contemporaneidade} aos objetos de conhecimento, superar a \textbf{fragmentação} e promover o desenvolvimento de \textbf{uma} visão sistêmica ( BRASIL, 2019 ).
Os TCTs estão dispostos em seis macroáreas temáticas, articuladas pela Coordenação-Geral de Educação Ambiental e Temas Transversais da Educação Básica, no Ministério da Educação. - Meio ambiente – Aborda as relações entre o \textbf{homem} e o meio ambiente, as relações de consumo e a compreensão do desenvolvimento sustentável. - Economia – Insere a educação financeira e fiscal relacionada ao mundo do trabalho. - Saúde – Traz o bem-estar físico, mental e social vinculado à educação alimentar e nutricional. - Cidadania e civismo – Trata dos valores e concepções na vida familiar e social, dos direitos humanos, da criança, do adolescente e do idoso.
Além da educação para o trânsito. - Multiculturalismo – Aborda o respeito à diversidade cultural, às matrizes históricas e culturais brasileiras. - Ciência e tecnologia – Defende a inclusão do letramento tecnológico, \textbf{que} envolve o aproveitamento das ferramentas tecnológicas a favor do estudante e da sociedade.
A incorporação dos TCTs \textbf{que} afetam a vida humana em escala local, regional e global nas práticas pedagógicas escolares deve ser feita, preferencialmente, de forma transversal e integradora, e se \textbf{basear} na \textbf{problematização} da realidade e das situações de aprendizagem ; na integração das habilidades e competências curriculares à resolução de problemas ; na superação da concepção fragmentada do conhecimento para \textbf{uma} visão sistêmica ; na promoção de \textbf{um} processo educativo continuado e do conhecimento como \textbf{uma} construção coletiva ( BRASIL, 2019 ).
Para a efetivação do processo de ensino e aprendizagem nas escolas, seja para incorporar os TCTs ou trabalhar os conteúdos da base comum, é fundamental a definição do \textbf{método} didático \textbf{que} será aplicado para \textbf{sistematização} do processo de ensino.
Isto porque o \textbf{método} determina o tipo de aprendizagem pretendida e quais ações organizadas de forma sistemática deverão ser desenvolvidas para \textbf{que} essa aprendizagem se efetive do modo esperado ( MARANHÃO, 2014 ).
O \textbf{método} didático é o fio condutor do trabalho pedagógico da escola e da ação docente ( MARCHIORATO, 2014 ).
Sendo assim, precisa estar vinculado aos objetivos educacionais, ao papel social e específico da escola e à concepção de aprendizagem.
Nesse contexto, as Diretrizes Curriculares Nacionais para Educação Básica sustentam \textbf{que} a escola tem a função de disseminar o conhecimento socialmente produzido e acumulado pela \textbf{humanidade} ( BRASIL, 2013 ).
Corroborando as DCN, as Diretrizes Curriculares do Estado do Maranhão consideram \textbf{que} a função social da escola \textbf{diz} respeito à apropriação dos elementos culturais essenciais à compreensão mais elaborada e \textbf{sistematizada} da realidade física, cultural, social, econômica e política, de forma a propiciar a ampliação da visão de mundo dos \textbf{sujeitos}.
Conforme a BNCC, a área de Ciências da Natureza e suas Tecnologias deve garantir ao estudante o desenvolvimento das competências específicas : 1.
Analisar fenômenos naturais e processos tecnológicos com base nas relações entre matéria e energia, para propor ações individuais e coletivas \textbf{que} aperfeiçoem processos produtivos, minimizem impactos socioambientais e melhorem as condições de vida em âmbito local, regional e/ou global. 2.
Construir e utilizar interpretações sobre a dinâmica da Vida, da Terra e do Cosmos para elaborar argumentos, realizar previsões sobre o funcionamento e a evolução dos seres vivos e do Universo, e fundamentar decisões éticas e responsáveis. 3.
Analisar situações-problema e avaliar aplicações do conhecimento científico e tecnológico e suas implicações no mundo, utilizando procedimentos e linguagens próprios das Ciências da Natureza, para propor soluções \textbf{que} considerem demandas locais, regionais e/ou globais, e comunicar suas descobertas e conclusões a públicos variados, em diversos contextos e por meio de diferentes mídias e tecnologias digitais de informação e comunicação ( TDIC ).
Destaca -se, também, \textbf{que} : > [... ] cabe aos sistemas e redes de ensino, assim como às escolas, em suas respectivas esferas de autonomia e competência, incorporar aos currículos e às propostas pedagógicas a \textbf{abordagem} de temas \textbf{contemporâneos} \textbf{que} afetam a vida humana em escala local, regional e global, preferencialmente de forma transversal e integradora. [... ] ( BRASIL, 2018 ).
Esses \textbf{fundamentos} orientam a escolha do \textbf{método} \textbf{dialético}, \textbf{que} tem a prática social como eixo do trabalho pedagógico.
Nesse sentido, a prática social é o ponto de partida e de chegada do processo de ensino \textbf{que} tem como finalidade ampliar a compreensão sobre os elementos, nexos, \textbf{inter-relações}, contradições e \textbf{fundamentos} \textbf{que} constituem a realidade social ( MARANHÃO, 2014 ).
Essa opção \textbf{metodológica} esclarece “ o movimento do conhecimento como a passagem do empírico ao concreto, pela mediação do abstrato.
Ou a passagem da síncrese à síntese, pela mediação da análise ” ( SAVIANI, 2011, p. 120 ).
Em outros termos, esse \textbf{método} consiste em partir da prática social, teorizar sobre ela e voltar à prática para transformá -la ( CORAZZA, 1991 apud GASPARIN, 2012 ).
Analisando o \textbf{método} \textbf{dialético} ( prática-teoria-prática ) à luz da \textbf{teoria} de Vygotsky, suas três fases coincidem com \textbf{um} processo \textbf{que} parte do nível de desenvolvimento real dos estudantes, trabalha na zona de desenvolvimento proximal, para chegar a \textbf{um} novo nível de desenvolvimento real.
Estes três momentos se desdobram nas cinco etapas do \textbf{método} \textbf{dialético} : prática social inicial, \textbf{problematização}, instrumentalização, catarse e prática social final.
Nesse \textbf{método}, a prática social inicial ( conhecimentos prévios dos estudantes ) é o ponto de partida do processo de ensino e aprendizagem.
Em seguida, a \textbf{problematização} permite selecionar e discutir os principais problemas dessa prática inicial, \textbf{que} serão respondidos na instrumentalização do conhecimento científico ( mediação pedagógica ).
A catarse é o momento de incorporação do conhecimento sintético, a verdadeira apropriação do saber pelo estudante, \textbf{que} retornará à prática social, agora, como prática social final, \textbf{uma} ação transformadora ( GASPARIN, 2012 ).
Nessa perspectiva \textbf{metodológica}, o estudante apreende novos conhecimentos a partir do confronto mental dos seus conhecimentos prévios com o conhecimento científico apresentado durante a prática pedagógica ( GASPARIN, 2012 ).
Assim, a aprendizagem se dá num processo contínuo de elaboração de relações entre os conhecimentos anteriores do estudante e as novas informações \textbf{que} são disponibilizadas no processo de ensino ( ABIB, 2010 ).
A aplicação do \textbf{método} \textbf{dialético} \textbf{exige} \textbf{uma} concepção de avaliação da aprendizagem a ser compreendida em três dimensões : avaliação inicial, avaliação processual e avaliação de resultado.
Cada dimensão da avaliação tem funções e instrumentos específicos, de acordo com a situação de aprendizagem.
Juntas, essas dimensões abrangem todos os momentos da prática pedagógica ( MARANHÃO, 2014 ).
A avaliação inicial é diagnóstica e deve ser feita no planejamento pedagógico, visando atender às necessidades de aprendizagens dos estudantes.
A avaliação processual deve ser formativa, possibilitando o acompanhamento dos avanços e das dificuldades dos estudantes durante a \textbf{problematização}, instrumentalização e catarse, com vistas às intervenções pedagógicas \textbf{que} melhorem o aprendizado.
A avaliação de resultado ocorre na prática social final e pretende verificar mudanças de atitude do estudante em relação às problemáticas estudadas ( MARANHÃO, 2014 ).
Acerca dos instrumentos avaliativos, a Lei de Diretrizes e Bases da Educação indica \textbf{que}, para a verificação do rendimento escolar, as avaliações deverão priorizar os aspectos qualitativos sobre os quantitativos.
Outra indicação é o uso de avaliações formativas de processo ou de resultado \textbf{que} levem em conta os contextos e as condições de aprendizagem, com o objetivo de melhorar o desempenho da escola, dos professores e dos estudantes ( BRASIL, 2017 ).
Considerando a diversidade de ações \textbf{didático-pedagógicas} e de metodologias adotadas para promover o ensino de Ciências da Natureza e suas Tecnologias, os instrumentos avaliativos deverão ser criteriosamente definidos para contemplar a variedade de situações de aprendizagem e, ao mesmo tempo, promover o melhor acompanhamento possível do aprendizado, de modo a assegurar a avaliação integral dos estudantes, numa perspectiva diagnóstica e formativa.
A definição e o uso dos instrumentos de avaliação precisam estar condicionados e adequados ao contexto e às condições de aprendizagem, aos objetivos e critérios de avaliação específicos da área e às competências e habilidades \textbf{que} se deseja avaliar.
Por exemplo, em atividades colaborativas e investigativas \textbf{que} envolvam a elaboração de modelos, protótipos e propostas de intervenção, o professor pode lançar mão de múltiplos instrumentos avaliativos, como pré-testes, pesquisas, questionários, relatórios ou trabalhos escritos, seminários, \textbf{exposições}, atividades de campo, avaliação por pares, rubricas, diários de registro, portfólio e autoavaliação. \textbf{Logo}, o ato de avaliar não pode ser compreendido como \textbf{simples} atribuição de nota por possível conhecimento obtido pelo estudante, mas como \textbf{um} processo fundamental, \textbf{que} deve ocorrer de forma contínua, possibilitando ao professor acompanhar os avanços e as dificuldades dos estudantes ao longo do processo, bem como promover, no tempo em \textbf{que} as dificuldades ocorrem, as intervenções pedagógicas necessárias, a fim de garantir o aprendizado esperado.
Os procedimentos \textbf{metodológicos} podem ser por desenvolvimento de projetos, estudo do meio, jogos, seminários, rodas de conversas com debates para promover a inquietação sobre questões locais e globais, análise e discussão de textos \textbf{que} tratam de questões como saúde e bem-estar físico, mental e social, sugerir simuladores e a construção de protótipos na resolução de situações-problema \textbf{que} envolvam sustentabilidade da vida no planeta.

% matched lemmas: abordagem, argumentar, atual, basear, compromisso, conceitual, contemporaneidade, contemporâneo, contextualização, corresponsável, crítico-reflexivo, despertar, dialético, didático-pedagógico, dizer, exigir, exposição, focar, fragmentação, frisar, fundamento, homem, humanidade, implicar, imprescindível, influenciar, influência, integralidade, inter-relacionar, inter-relação, interdisciplinaridade, logo, metodológico, método, preciso, problematização, que, simples, sistematizar, sistematização, sujeito, século, sócio-histórico, teoria, teórico, transdisciplinaridade, um
\end{textsample}
